% Discussion, limitations, conclusion. Drafted 2026-09-05; prose revised 2026-09-05.
\section{Discussion}\label{sec:discussion}

\paragraph{What the certificate says and what the decoder does are decoupled}
Nothing in this audit contradicts Theorem~3.1 of He et al.: the implementation does what it claims on every decode step we checked. The guarantee, as stated, describes the mechanism poorly in both directions. It promises too little: Eq.~\eqref{eq:knaf} bounds an expectation, so single outputs exceed $K$ (Section~\ref{sec:tails}), and its consequence for the rights-holder's event, Eq.~\eqref{eq:cap}, is vacuous for every CopyBench passage at the budgets of the released runs, under either anchor we tried (Section~\ref{sec:certificate}). It also understates what the decoder does: at those budgets, against risky models that memorised the passages, the banked budget binds and single-query reproduction stays below $10\%$ (Sections~\ref{sec:attack-results} and~\ref{sec:natural}). The protection that exists comes from the anchor's surprisal of the passage exceeding the total budget by a margin the certificate does not express, and Section~\ref{sec:utility} prices the consequence: the budgets at which the mechanism preserves utility are the budgets at which it does nothing, and the budgets at which the guarantee is informative cost nearly as much as serving the anchor alone.

\paragraph{Composition, the prefix debt, and what a budget cannot separate}
Cohen's theorem~\cite{cohen2025blameless} predicts that a per-prompt guarantee cannot bound what many prompts reveal, and the oracle-prefix attack confirms it: composition doubles the single query's take at $k=5$ and $k=10$ and recovers $86\%$ of a passage at $k=20$, every query compliant. The prefix debt, introduced to stop banking across a copied prompt, does more against extraction on $k \in [3,10]$ than the budget does and nothing at all at $k=20$ (Section~\ref{sec:prefixdebt}); it taxes exactly the prompts an extraction attack must use, so a guarantee that accounted for it would have to say what it assumes about the prompt. The natural repair for composition is the odometer of differential privacy, and Proposition~\ref{prop:sep} says what it can and cannot buy. It can bound reconstruction. It cannot do so selectively, because ordinary use and reconstruction are charged in the same additive nats at rates whose ratio $k$ fixes rather than the deployer: at $k=20$ the budget that holds the strongest adversary to a tenth of a work admits $0.61$ of one ordinary answer. Escaping that needs an accounting that is not the decoder's own --- a semantic detector, or a budget attached to the work rather than the user.

\paragraph{Three repairs, and what they are worth}
Charging the realised log-ratio turns the certificate into the $\Delta_{\max}$ form, so the same decoder, budgets and banking rule guarantee for every output what the KL decoder delivers on average, for a small utility loss concentrated where the KL decoder pays nothing (Sections~\ref{sec:pathwise} and~\ref{sec:utility}). Concentration is the weaker option and reaches the wrong place: Proposition~\ref{prop:conc} is informative at every deployed budget yet silent about the outputs that already exceed $K$, which are the ones being asked about. Neither form restores meaning to the certificate for short passages at $k \ge 3$, because both bound reproduction through the anchor's surprisal, which the budget exceeds and a better-read anchor lowers rather than raises.

\paragraph{A certificate card}
$K$ alone does not tell a rights-holder what they are protected against, and the missing information is cheap to publish. Following the epsilon registry proposed for differential privacy~\cite{dwork2019expose}, a deployment should publish next to $K$ the anchor and the decoding settings the certificate is relative to, the prefix-debt rule, $k$ and $T_{\max}$, the active-step and identical-output fractions at the deployed $k$, the vacuity fraction and per-token anchor surprisal on a reference set of protected passages, the concentration bound of Proposition~\ref{prop:conc} or the $\Delta_{\max}$ form if enforced, any per-user budget and bank cap, and for works the risky model is known to have memorised, the reconstruction C7 measures. A rights-holder can compute the third and fourth of those from the anchor alone.

\section{Limitations}\label{sec:limitations}
We audit one implementation and two risky models: a fine-tuned 8B model, whose memorisation is the adversary's choice, and the 70B base model the mechanism's authors evaluate, on the two CopyBench books it is known to have memorised. Only TinyComma can be fused with a Llama-3 risky model --- the 7B Common Pile anchor has a $64{,}000$-token vocabulary against Llama-3's $128{,}256$, and the decoder needs one shared vocabulary --- so the attacks use a single anchor and the second appears only where surprisal is computed. The second mechanism is our reimplementation of the published rule, not the authors' code. $S(x)$ is measured to $1{,}024$ tokens, the anchor's context, so longer works are extrapolated rather than observed, and near-verbatim recall uses a $20$-word span because the passages are short. Utility is judged by one instruction-tuned model on the benchmark's own prompt classes rather than a deployment's traffic, and one judge is a single point of failure even with the null of Section~\ref{sec:utility}. We make no claim about legal thresholds for infringement and do not evaluate non-literal reproduction.
